\documentclass[10pt,a4paper]{amsart}
\usepackage{amsmath,amssymb,amsthm,graphicx}
% Retained minimal preamble for the Stacks Project chapter flat.tex (More on Flatness).
% Verbatim excerpt of the chapter preamble: Xy-pic setup (native preamble lines 14-18) and the native
% theorem declarations with the mathematical macros (native preamble lines 154-217). Copied mechanically;
% nothing added, removed or redefined. The wrapper class and the amsmath/amssymb/amsthm/graphicx base
% packages are supplied by the assembler, so the class/externaldocument part of the native preamble is omitted.
% For commutative diagrams we use Xy-pic
\usepackage[all]{xy}

% We use 2cell for 2-commutative diagrams.
\xyoption{2cell}
% Theorem environments.
%
\theoremstyle{plain}
\newtheorem{theorem}[subsection]{Theorem}
\newtheorem{proposition}[subsection]{Proposition}
\newtheorem{lemma}[subsection]{Lemma}

\theoremstyle{definition}
\newtheorem{definition}[subsection]{Definition}
\newtheorem{example}[subsection]{Example}
\newtheorem{exercise}[subsection]{Exercise}
\newtheorem{situation}[subsection]{Situation}

\theoremstyle{remark}
\newtheorem{remark}[subsection]{Remark}
\newtheorem{remarks}[subsection]{Remarks}

\numberwithin{equation}{subsection}

% Macros
%
\def\lim{\mathop{\mathrm{lim}}\nolimits}
\def\colim{\mathop{\mathrm{colim}}\nolimits}
\def\Spec{\mathop{\mathrm{Spec}}}
\def\Hom{\mathop{\mathrm{Hom}}\nolimits}
\def\Ext{\mathop{\mathrm{Ext}}\nolimits}
\def\SheafHom{\mathop{\mathcal{H}\!\mathit{om}}\nolimits}
\def\SheafExt{\mathop{\mathcal{E}\!\mathit{xt}}\nolimits}
\def\Sch{\mathit{Sch}}
\def\Mor{\mathop{\mathrm{Mor}}\nolimits}
\def\Ob{\mathop{\mathrm{Ob}}\nolimits}
\def\Sh{\mathop{\mathit{Sh}}\nolimits}
\def\NL{\mathop{N\!L}\nolimits}
\def\CH{\mathop{\mathrm{CH}}\nolimits}
\def\proetale{{pro\text{-}\acute{e}tale}}
\def\etale{{\acute{e}tale}}
\def\QCoh{\mathit{QCoh}}
\def\Ker{\mathop{\mathrm{Ker}}}
\def\Im{\mathop{\mathrm{Im}}}
\def\Coker{\mathop{\mathrm{Coker}}}
\def\Coim{\mathop{\mathrm{Coim}}}

% Boxtimes
%
\DeclareMathSymbol{\boxtimes}{\mathbin}{AMSa}{"02}

%
% Macros for moduli stacks/spaces
%
\def\QCohstack{\mathcal{QC}\!\mathit{oh}}
\def\Cohstack{\mathcal{C}\!\mathit{oh}}
\def\Spacesstack{\mathcal{S}\!\mathit{paces}}
\def\Quotfunctor{\mathrm{Quot}}
\def\Hilbfunctor{\mathrm{Hilb}}
\def\Curvesstack{\mathcal{C}\!\mathit{urves}}
\def\Polarizedstack{\mathcal{P}\!\mathit{olarized}}
\def\Complexesstack{\mathcal{C}\!\mathit{omplexes}}
% \Pic is the operator that assigns to X its picard group, usage \Pic(X)
% \Picardstack_{X/B} denotes the Picard stack of X over B
% \Picardfunctor_{X/B} denotes the Picard functor of X over B
\def\Pic{\mathop{\mathrm{Pic}}\nolimits}
\def\Picardstack{\mathcal{P}\!\mathit{ic}}
\def\Picardfunctor{\mathrm{Pic}}
\def\Deformationcategory{\mathcal{D}\!\mathit{ef}}

% --- appended 2026-10-12 (Chief): the project own definitions of the reference/slogan/history
% display environments, copied mechanically from the shared Stacks preamble lines 8-12
% (\usepackage{verbatim} plus the three \newenvironment lines, which the chapter preamble keeps
% verbatim from the project preamble). Nothing added, removed or redefined: these environments
% carry no mathematical content and the project itself typesets them as comments. Without these
% three lines a counted statement whose own text uses them fails with Environment slogan undefined.
\usepackage{verbatim}
\newenvironment{reference}{\comment}{\endcomment}
\newenvironment{slogan}{\comment}{\endcomment}
\newenvironment{history}{\comment}{\endcomment}
\usepackage{cleveref}
\crefname{originallemma}{original lemma}{original lemmas}
\Crefname{originallemma}{Original lemma}{Original lemmas}
\makeatletter
\newlabel{lemma-complete-devissage-flat-finitely-presented-module}{{05IB}{source}{}{sourcepointer-1}{}}
\newlabel{lemma-complete-devissage-flat-finitely-presented-module@cref}{{[originallemma][0][]05IB}{[1][0][]source}{}{}{}}
\makeatother
\crefname{originallemma}{original lemma}{original lemmas}
\Crefname{originallemma}{Original lemma}{Original lemmas}
\makeatletter
\newlabel{lemma-etale-flat-up-down}{{05B9}{source}{}{sourcepointer-2}{}}
\newlabel{lemma-etale-flat-up-down@cref}{{[originallemma][0][]05B9}{[1][0][]source}{}{}{}}
\makeatother
\crefname{originaltheorem}{original theorem}{original theorems}
\Crefname{originaltheorem}{Original theorem}{Original theorems}
\makeatletter
\newlabel{more-morphisms-theorem-openness-flatness}{{0399}{source}{}{sourcepointer-3}{}}
\newlabel{more-morphisms-theorem-openness-flatness@cref}{{[originaltheorem][0][]0399}{[1][0][]source}{}{}{}}
\makeatother
\crefname{originallemma}{original lemma}{original lemmas}
\Crefname{originallemma}{Original lemma}{Original lemmas}
\makeatletter
\newlabel{morphisms-lemma-flat-permanence}{{02JZ}{source}{}{sourcepointer-4}{}}
\newlabel{morphisms-lemma-flat-permanence@cref}{{[originallemma][0][]02JZ}{[1][0][]source}{}{}{}}
\makeatother
\crefname{originalproposition}{original proposition}{original propositions}
\Crefname{originalproposition}{Original proposition}{Original propositions}
\makeatletter
\newlabel{proposition-existence-complete-at-x}{{05HR}{source}{}{sourcepointer-5}{}}
\newlabel{proposition-existence-complete-at-x@cref}{{[originalproposition][0][]05HR}{[1][0][]source}{}{}{}}
\makeatother
\crefname{originalremark}{original remark}{original remarks}
\Crefname{originalremark}{Original remark}{Original remarks}
\makeatletter
\newlabel{remark-same-notion}{{05I0}{source}{}{sourcepointer-6}{}}
\newlabel{remark-same-notion@cref}{{[originalremark][0][]05I0}{[1][0][]source}{}{}{}}
\makeatother
\title{Projective global sections after an elementary etale neighbourhood for a finitely presented module flat at a point}
\author{The Stacks Project Authors}
\date{stacks-project:flat.tex:local-snapshot-20261011; GNU FDL 1.2 or later}
\begin{document}
\maketitle
\noindent\textit{Editorial note.} All delivered mathematical text is the author's own, copied from the pinned Stacks Project chapter file flat.tex; no statement or proof step was rewritten, reordered or removed.\par\medskip
\noindent\textit{Editorial note.} Preamble provenance: the chapter's own retained minimal preamble (sources/stacks-flat/retained-preamble.tex, lines 1-85) is a verbatim excerpt of the native Xy-pic setup and of the native theorem declarations and mathematical macros, followed by the project's own reference/slogan/history display-environment definitions copied mechanically from the shared preamble; the wrapper class and the amsmath/amssymb/amsthm/graphicx base packages are supplied by the assembler, and extra\_packages is left empty because everything needed is already in the retained preamble. No macro was invented or redefined.\par\medskip
\noindent\textit{Editorial note.} Same-chapter pointers are located at their \textbackslash{}label lines inside the pinned chapter file; cross-chapter pointers are located in the project's own tag list staged in the snapshot (sources/stacks-chow/tags.txt), at the verbatim lines carrying their canonical tags.\par\medskip
\noindent\textit{Allowed standard prerequisites (editorial).}\begin{itemize}
\item Openness of flatness. More on Morphisms, label more-morphisms-theorem-openness-flatness, invoked at proof line 3312. Lets the module be assumed flat over the whole of an open neighbourhood of the point, which is the first reduction.
\item Existence of a complete devissage after an elementary etale neighbourhood. Same chapter flat.tex, label proposition-existence-complete-at-x, invoked at proof line 3315. Produces the diagram of pointed schemes whose horizontal arrows are elementary etale neighbourhoods and over which the pulled back module has a complete devissage at the lifted point.
\item Permanence of flatness, the up and down lemma for an etale cover, the identification of the affine description, and the consequences of a complete devissage. Morphisms, label morphisms-lemma-flat-permanence (proof line 3319); Same chapter flat.tex, labels lemma-etale-flat-up-down (proof line 3321, claimed as S14-algebra-etale-flat-up-down), remark-same-notion (proof line 3323) and lemma-complete-devissage-flat-finitely-presented-module (proof line 3332, claimed as S14-algebra-devissage-fp-flat). They keep flatness after the etale base change, transfer it to the new base, translate the devissage into the affine-algebra statement, and then conclude projectivity of the module over the new base after replacing it by a standard open neighbourhood.
\end{itemize}\medskip
\ifdefined\hypertarget\hypertarget{sourcepointer-1}{}\fi
\noindent\textit{Source reference.} Original lemma 05IB: sources/stacks-flat/pinned-original-source.tex, source lines 3229--3229. Navigation pointer for \textbackslash{}ref\{lemma-complete-devissage-flat-finitely-presented-module\} used inside the counted span (line 3332); the argument that consumes it is retained in the span.\par\medskip
\ifdefined\hypertarget\hypertarget{sourcepointer-2}{}\fi
\noindent\textit{Source reference.} Original lemma 05B9: sources/stacks-flat/pinned-original-source.tex, source lines 124--124. Navigation pointer for \textbackslash{}ref\{lemma-etale-flat-up-down\} used inside the counted span (line 3321); the argument that consumes it is retained in the span.\par\medskip
\ifdefined\hypertarget\hypertarget{sourcepointer-3}{}\fi
\noindent\textit{Source reference.} Original theorem 0399: sources/stacks-chow/tags.txt, source lines 4264--4264. Navigation pointer for \textbackslash{}ref\{more-morphisms-theorem-openness-flatness\} used inside the counted span (line 3312); the argument that consumes it is retained in the span.\par\medskip
\ifdefined\hypertarget\hypertarget{sourcepointer-4}{}\fi
\noindent\textit{Source reference.} Original lemma 02JZ: sources/stacks-chow/tags.txt, source lines 3354--3354. Navigation pointer for \textbackslash{}ref\{morphisms-lemma-flat-permanence\} used inside the counted span (line 3319); the argument that consumes it is retained in the span.\par\medskip
\ifdefined\hypertarget\hypertarget{sourcepointer-5}{}\fi
\noindent\textit{Source reference.} Original proposition 05HR: sources/stacks-flat/pinned-original-source.tex, source lines 1429--1429. Navigation pointer for \textbackslash{}ref\{proposition-existence-complete-at-x\} used inside the counted span (line 3315); the argument that consumes it is retained in the span.\par\medskip
\ifdefined\hypertarget\hypertarget{sourcepointer-6}{}\fi
\noindent\textit{Source reference.} Original remark 05I0: sources/stacks-flat/pinned-original-source.tex, source lines 1731--1731. Navigation pointer for \textbackslash{}ref\{remark-same-notion\} used inside the counted span (line 3323); the argument that consumes it is retained in the span.\par\medskip
\begin{proposition}
\label{proposition-finite-presentation-flat-at-point}
Let $f : X \to S$ be a morphism of schemes.
Let $\mathcal{F}$ be a quasi-coherent sheaf on $X$.
Let $x \in X$ with image $s \in S$.
Assume that
\begin{enumerate}
\item $f$ is locally of finite presentation,
\item $\mathcal{F}$ is of finite presentation, and
\item $\mathcal{F}$ is flat at $x$ over $S$.
\end{enumerate}
Then there exists a commutative diagram of pointed schemes
$$
\xymatrix{
(X, x) \ar[d] & (X', x') \ar[l]^g \ar[d] \\
(S, s) & (S', s') \ar[l]
}
$$
whose horizontal arrows are elementary \'etale neighbourhoods
such that $X'$, $S'$ are affine and such that
$\Gamma(X', g^*\mathcal{F})$ is a projective
$\Gamma(S', \mathcal{O}_{S'})$-module.
\end{proposition}

\begin{proof}
By openness of flatness, see
More on Morphisms, Theorem \ref{more-morphisms-theorem-openness-flatness}
we may replace $X$ by an open neighbourhood of $x$ and assume that
$\mathcal{F}$ is flat over $S$. Next, we apply
Proposition \ref{proposition-existence-complete-at-x}
to find a diagram as in the statement of the proposition such
that $g^*\mathcal{F}/X'/S'$ has a complete d\'evissage over $s'$.
(In particular $S'$ and $X'$ are affine.) By
Morphisms, Lemma \ref{morphisms-lemma-flat-permanence}
we see that $g^*\mathcal{F}$ is flat over $S$ and by
Lemma \ref{lemma-etale-flat-up-down}
we see that it is flat over $S'$. Via
Remark \ref{remark-same-notion}
we deduce that
$$
\Gamma(X', g^*\mathcal{F})/
\Gamma(X', \mathcal{O}_{X'})/
\Gamma(S', \mathcal{O}_{S'})
$$
has a complete d\'evissage over the prime of $\Gamma(S', \mathcal{O}_{S'})$
corresponding to $s'$. Thus
Lemma \ref{lemma-complete-devissage-flat-finitely-presented-module}
implies that the result of the proposition holds after replacing
$S'$ by a standard open neighbourhood of $s'$.
\end{proof}

\end{document}
